\section{The same native generic coordinate ratios across original projective charts}
\label{sec:projective-chart-generic-transitions}

\begin{theorem}[Exact original generic coordinate transitions]
\label{thm:projective-chart-generic-transitions}
\lean{EnvelopeProperty.ProjectiveChartGenericTransitions.coordinate_transition, EnvelopeProperty.ProjectiveChartGenericTransitions.coordinate_pivot_ne_zero, EnvelopeProperty.ProjectiveChartGenericTransitions.coordinate_div}
\leanok
\uses{thm:fc-original-homogeneous-form-section,thm:projective-chart-function-field-generation}
\source{bj22-trivval:sec:tc}
For an actual morphism $f:X\to\mathbf P^m_k$ from an integral original
scheme and any two coordinate opens with nonempty original inverse
images, let $c_{i,a}$ be the literal original generic coefficient
$X_a/X_i$ in the native $K(X)$. Then
\[
 c_{i,a}=c_{i,j}c_{j,a},\qquad c_{i,j}\ne0,
 \qquad c_{i,a}/c_{i,j}=c_{j,a}.
\]
No original closed immersion, independent field comparison or coordinate
transition premise is required.
\end{theorem}
\begin{proof}
The original generic point belongs to both nonempty original opens and
therefore to their actual overlap. Restricting a whole original ambient
section before taking its pulled native generic germ gives the same
germ as pulling the original section on the larger open: the actual
native stalk-map/germ square and sheaf restriction give this identity.

The already proved whole homogeneous-section transition identity on
the original overlap says that the coefficient of $X_a$ in chart $i$
is the original transition unit $X_j/X_i$ times its coefficient in
chart $j$. Apply the actual original pulled generic-germ algebra map
on this same overlap. The restriction identity identifies its three
terms with $c_{i,a}$, $c_{i,j}$ and $c_{j,a}$, giving the first formula.
Taking $a=i$ and using the actual original pivot coefficient one makes
$c_{i,j}c_{j,i}=1$, proving nonvanishing. Division in the same original
native field gives the final formula. Every coordinate and field is
the original one throughout.
\end{proof}
